% AUTORIDAD DE FUENTE:
% expose_VIII_body.tex, líneas 2316--2361; índices de escaneo 315--317;
% folios de fuente 304--306. Los dos E/E_eta de la misma disposición
% francesa se conservan reversiblemente mediante \sisi.

\medskip
\noindent 7.3.2. De la compatibilidad (7.3.1.1) se sigue que, si
$E_\eta$ es una extensión de $P_\eta$ por $\underline G_{m\eta}$,
existe un mayor subesquema en grupos $V$ de $P$ tal que la restricción
de $E_\eta$ a $V_\eta$ se prolonga a una extensión de $V$ por
$\underline G_{mS}$. Es el subgrupo abierto $V$ de $P$ que contiene
$P_\eta$ y que viene definido por la condición
\begin{equation}\tag{7.3.2.1}
 V_0/V_0^{\circ}=\ker d(E_\eta).
\end{equation}
Si $u:P'\to P$ es un morfismo de esquemas en grupos lisos sobre $S$,
la extensión $E'_\eta$, obtenida por imagen inversa de $E_\eta$ mediante
$u_\eta$, se prolonga a una extensión de $P'$ por $\underline G_{mS}$
si y solo si
\begin{equation}\tag{7.3.2.2}
 u(P')\subset V.
\end{equation}

%%%%%%%%%%  scan idx 316  |  folio 305  %%%%%%%%%%

En el caso de una biextensión $E\sisi{}{_\eta}$ de
$(P_\eta,Q_\eta)$ por $\underline G_{m\eta}$, existe asimismo un mayor
subesquema en grupos $V$ de $P$ tal que la restricción de
$E\sisi{}{_\eta}$ a $(V_\eta,Q_\eta)$ se prolonga a una biextensión de
$(V,Q)$ por $\underline G_{mS}$. Es el subgrupo abierto de $P$ que
contiene la fibra genérica $P_\eta$ y está caracterizado por que
$V_0/V_0^{\circ}$ sea el anulador izquierdo en $\Phi$ del
emparejamiento (7.1.2), es decir, el núcleo del homomorfismo
correspondiente de $\Phi$ a
$\underline{\Hom}(\Psi,(\mathbb Q/\mathbb Z)_k)$:
\begin{equation}\tag{7.3.2.3}
 V_0/V_0^{\circ}
 =\ker\!\left(
   \Phi\xrightarrow{\ d(E_\eta)\ }
   \underline{\Hom}(\Psi,(\mathbb Q/\mathbb Z)_k)
 \right).
\end{equation}
Si $u:P'\to P$ es un morfismo de esquemas en grupos lisos sobre $S$,
la biextensión de $(P'_\eta,Q_\eta)$ por $\underline G_{m\eta}$
obtenida por imagen inversa de $E_\eta$ mediante $u_\eta$ se prolonga
a una biextensión de $(P',Q)$ por $\underline G_{mS}$ si y solo si se
cumple (7.3.2.2). Se obtienen resultados simétricos intercambiando los
papeles de $P$ y $Q$.

%%%%%%%%%%  scan idx 317  |  folio 306  %%%%%%%%%%

Naturalmente, en general no existe una mayor \underline{pareja} de
subgrupos $P'\to P$, $Q'\to Q$ para la cual la restricción
$E'_\eta$ de $E_\eta$ a una biextensión de
$(P'_\eta,Q'_\eta)$ por $\underline G_{m\eta}$ se prolongue a una
biextensión de $(P',Q')$ por $\underline G_{mS}$. Solo puede decirse
que, para una pareja dada $(P',Q')$, o más generalmente para una pareja
dada de morfismos $P'\to P$, $Q'\to Q$ de esquemas en grupos lisos, la
prolongación buscada existe si y solo si la restricción de la forma
$d(E_\eta)$ a
\[
 P'_0/{P'_0}^{\circ}\otimes Q'_0/{Q'_0}^{\circ}
\]
es nula.

\medskip
\noindent 7.3.3. \underline{Simetría}. Sea $E_\eta$ una biextensión de
$(P_\eta,Q_\eta)$ por $\underline G_{m\eta}$, y consideremos la
biextensión transpuesta ${}^{s}E_\eta$ (VII 2.7), que es una
biextensión de $(Q_\eta,P_\eta)$ por $\underline G_{m\eta}$. El
diagrama siguiente es conmutativo:
\begin{equation}\tag{7.3.3.1}
\begin{tikzcd}[row sep=large,column sep=huge]
 \Phi\otimes\Psi
   \arrow[dd,"\mathrm{sym}"']
   \arrow[rd,"d(E_\eta)"] & \\
 & (\mathbb Q/\mathbb Z)_k \\
 \Psi\otimes\Phi
   \arrow[ru,"d({}^{s}E_\eta)"'] &
\end{tikzcd}
\qquad ,
\end{equation}
donde la flecha vertical es la simetría canónica. Compárese esta
afirmación con la de \underline{antisimetría} de 2.2.10. Se sigue, en
particular, que si $P=Q$ y $E_\eta$ es una biextensión
\underline{simétrica} de $(P_\eta,Q_\eta)$ por
$\underline G_{m\eta}$, es decir, si $E_\eta$ es isomorfa a
${}^{s}E_\eta$, entonces la forma bilineal $d(E_\eta)$ sobre
$\Phi\times\Phi$ con valores en $(\mathbb Q/\mathbb Z)_k$ es
\underline{simétrica} (¡y no alternada, como en 2.2.12!). Si es nula,
es decir, si $E_\eta$ se prolonga a una biextensión $E$ de $(P,P)$ por
$\underline G_{mS}$, esta última es necesariamente también simétrica:
por 7.1(a), todo isomorfismo ${}^{s}E_\eta\to E_\eta$ se prolonga de
manera única a un isomorfismo ${}^{s}E\to E$.

%%%%%%%%%%  scan idx 318  |  folio 307  %%%%%%%%%%
